\chapter{A finite-record hierarchy of classical realization size}
\label{ch:finite-hierarchy}

The three-state occupation bound can exclude a small classical explanation,
but one scalar contrast cannot establish an unbounded hierarchy. Retaining
several observation times changes the mathematical object. A single prepared
state, a single binary detector, and a common repeated evolution produce a
matrix with a positive factorization through the hidden state space. This
chapter turns a robust obstruction to that factorization into a finite-record
experiment. The dimension is the total number of states of the proposed
classical Markov explanation, including any hidden controller or clock.

The foundations are established: positive realization and nonnegative-rank
bounds, Sperner's antichain theorem, and quantum dimension witnesses already
connect small quantum systems to larger classical realizations~\cite{rfGillis2020,rfShahandeh2024}. In particular,
the zero-noise sine matrix used below has appeared in prior work, including a
stronger six-factor lower bound for its seven-by-seven instance~\cite{rfHeinosaari2024}. The scoped
addition developed here is a nonzero-noise certificate, its realization through
one repeated channel and finitely many observation times, and an exactly checked
finite-record statistical protocol. Historical priority of that combination is
not established by this internal derivation.

\section{The operational model}
Let $P$ be any $n\times n$ stochastic matrix, $\alpha$ a probability row, and
$f\in[0,1]^n$ a fixed stochastic binary readout. Independent reset experiments
observe a binary outcome after $k$ applications of the same channel, with mean
\begin{equation}
 p_k=\alpha P^k f.                                      \label{eq:hier-pk}
\end{equation}
A continuous-time homogeneous chain is included by $P=e^{\Delta Q}$. The
result also covers repeated stochastic interventions, and independent,
identically distributed nonnegative timing increments with
$P=\mathbb E[e^{QZ}]$. It does not require a generator, stationarity, reversibility,
known rates, pure preparations, or a noiseless detector.

Fix $N\ge2$ and measure $p_1,\ldots,p_{2N-1}$. Form the matrix
\begin{equation}
 M_{ij}=p_{N+i-j},\qquad 0\le i,j<N.
                                                        \label{eq:hier-matrix}
\end{equation}
It factors as $M=AB$, where $A_i=\alpha P^i$ and $B_j=P^{N-j}f$.
Every row of $A$ is a probability vector and every entry of $B$ lies in $[0,1]$.
The reversal of the column index is only a permutation of a finite Hankel
matrix. There are $2N-1$ distinct measurement settings, not $N^2$ distinct
experiments. The diagonal entries all equal $p_N$.

\section{A robust antichain certificate}
\begin{theorem}[Finite-noise positive-factorization obstruction]
\label{thm:robust-antichain}
Suppose an $N\times N$ matrix admits $M=AB$ with $n$ columns in $A$,
probability rows in $A$, and $0\le B_{kj}\le1$. If
\begin{equation}
 M_{ii}\le\epsilon,\qquad M_{ij}\ge\beta\quad(i\ne j),
 \qquad
 \beta>\tau+(1-\tau)\frac{\epsilon}{\tau}
                                                        \label{eq:hier-margin}
\end{equation}
for some $0<\tau<1$, then
\begin{equation}
 N\le {n\choose\lfloor n/2\rfloor}.                     \label{eq:hier-sperner}
\end{equation}
For $0<\epsilon<1$, minimizing the right side of
\eqref{eq:hier-margin} over $\tau$ gives the sufficient condition
$\beta>2\sqrt\epsilon-\epsilon$. For $\epsilon=0$, any $\beta>0$
suffices.
\end{theorem}
\begin{proof}
Write $a_i$ for row $i$ of $A$, and define
$S_i=\{k:B_{ki}<\tau\}$. The diagonal bound gives
$a_i(S_i^c)\le\epsilon/\tau$. If $S_i\subseteq S_j$, then
\[
 M_{ij}\le\tau a_i(S_j)+a_i(S_j^c)
 \le\tau+(1-\tau)\epsilon/\tau,
\]
contradicting \eqref{eq:hier-margin}. Thus the $N$ subsets $S_i$ are pairwise
incomparable. Sperner's theorem bounds their number by the largest binomial
coefficient. The minimum of $\tau+\epsilon/\tau-\epsilon$ is obtained at
$\tau=\sqrt\epsilon$. When $\epsilon=0$, take $S_i=\{k:B_{ki}=0\}$.
The diagonal condition puts all of $a_i$ on $S_i$, and inclusion
$S_i\subseteq S_j$ would give $M_{ij}=0$.
\end{proof}

This certificate supplies a sufficient lower bound, not the nonnegative rank
itself and not the minimum dimension of a positive dynamical realization.
Those quantities need not coincide. The antichain cardinality bound is sharp
for unrestricted stochastic factors: use all subsets of size $\lfloor n/2\rfloor$,
let $a_i$ be uniform on subset $S_i$, and let $B_{kj}$ indicate $k\notin S_j$.
The diagonal is zero and every off-diagonal entry is positive. This construction
does not establish sharpness for a single-channel Hankel matrix.

\begin{corollary}[One-channel finite-time dimension hierarchy]
\label{cor:hierarchy}
If the records \eqref{eq:hier-pk} satisfy $p_N\le\epsilon$ and
$p_k\ge\beta$ for every $1\le k\le2N-1$, $k\ne N$, with
\eqref{eq:hier-margin}, their classical homogeneous Markov realization requires
at least
\[
 d_*(N)=\min\left\{n:{n\choose\lfloor n/2\rfloor}\ge N\right\}
\]
hidden states. In particular $N=7,11,21,71$ require respectively at least
$5,6,7,9$ states.
\end{corollary}

\section{An explicit qubit family}
Prepare a qubit in $|0\rangle$, use a fixed $|1\rangle$ detector, and evolve
with
\[
 \mathcal L(\rho)=-i\frac{\omega}{2}[X,\rho]
 +\gamma\bigl(\operatorname{tr}(\rho)I/2-\rho\bigr).
\]
Set $r=\omega/\gamma$ and $\Delta=2\pi/(N\omega)$. The measured means are
\begin{equation}
 p_k=\frac{1-e^{-2\pi k/(rN)}\cos(2\pi k/N)}2.
                                                        \label{eq:hier-qubit}
\end{equation}
For $N\ge4$, their diagonal value is
$\epsilon_r=(1-e^{-2\pi/r})/2$, and all the off-diagonal values are at least
$\sin^2(\pi/N)$. If the cosine is nonnegative, damping increases the probability
above its unitary value; if it is negative, the probability remains at least
$1/2\ge\sin^2(\pi/N)$. It follows that any fixed desired finite state lower
bound can be reached at a sufficiently large, but finite, coherence ratio.
This is a hierarchy of finite-resource experiments, not one finite experiment
that excludes all finite classical models.

The sufficient coherence condition can be written
\begin{equation}
 \epsilon_r < \bigl(1-\cos(\pi/N)\bigr)^2.
                                                        \label{eq:hier-coherence}
\end{equation}
For large $N$ its right side is of order $N^{-4}$. Since $N$ itself grows
roughly exponentially in the desired state count, this certificate is not an
efficient route to arbitrarily large dimension exclusions. Its transparent
finite-noise guarantee is the benefit; efficiency and optimality remain
separate questions. No relation between $r$ and thermodynamic energy is assumed.

The exact-arithmetic implementation gives the following demonstrations:
\begin{center}
\begin{tabular}{rrrr}
$N$ & observation settings & $\omega/\gamma$ & certified states\\
\hline
7 & 13 & $10^3$ & at least 5\\
11 & 21 & $10^4$ & at least 6\\
21 & 41 & $10^5$ & at least 7
\end{tabular}
\end{center}
They are calculated qubit probabilities. None is a laboratory observation.

\section{A finite-record test without a multiplicity penalty on type I error}
Choose a forbidden rectangular probability region
$p_N\le\epsilon_0$ and $p_k\ge\beta_0$ for $k\ne N$, satisfying
\eqref{eq:hier-margin}. For each setting collect an independent fixed number
$n_k$ of reset records, and let $S_k$ be the success count. Reject the proposed
classical dimension only when
\begin{equation}
 S_N\le c_N\quad\hbox{and}\quad S_k\ge c_k\quad(k\ne N).
                                                        \label{eq:hier-test}
\end{equation}
Calibrate each individual tail at its boundary to be at most $\alpha$.
Every model in the excluded classical class violates at least one of the
probability constraints. The probability of passing every count criterion is
therefore at most the rejection probability for that one violated constraint,
which is at most $\alpha$. A Bonferroni division of $\alpha$ among settings is
unnecessary. This is an intersection--union test, an established statistical
construction. For a specified alternative, the probability of failing any
criterion is bounded by the sum of its individual miss probabilities.

Here is a fully specified example. Set $N=7$, $r=1000$,
$\epsilon_0=3/500$, $\beta_0=3/20$, and $\tau=2/25$.
The classical comparison bound is $149/1000<\beta_0$.
Use the following counts and acceptance conditions:
\begin{center}
\begin{tabular}{rrrr}
Step $k$ & Records $n_k$ & Critical count $c_k$ & Requirement\\
\hline
1 &1900&313&$S_1\ge313$\\
2--5 &100 each&22&$S_k\ge22$\\
6 &1800&297&$S_6\ge297$\\
7 &10300&48&$S_7\le48$\\
8 &1800&297&$S_8\ge297$\\
9--12 &100 each&22&$S_k\ge22$\\
13 &1700&282&$S_{13}\ge282$
\end{tabular}
\end{center}
At $k=7$ require $S_7\le48$; elsewhere require $S_k\ge c_k$.
The total is $18\,300$ binary records and $126\,900$ channel applications.
The exact certificates imply
\begin{equation}
 \Pr_{n\le4}(\text{reject})<0.049048,
 \qquad
 \Pr_{\text{specified qubit}}(\text{do not reject})<0.015708.
                                                        \label{eq:hier-errors}
\end{equation}
The elapsed evolution time is $126900\Delta$, excluding preparation,
reset, readout, and controller overhead. An energy cost cannot be inferred
from this operation count without a physical implementation and calibration.
The record count is a checked sufficient design; no global optimality is claimed.

The probability intervals are certified with rational Machin-formula bounds
on $\pi$ and Taylor remainders. A binomial tail is bounded using its boundary
mass and a geometric series whose ratio bounds all successive masses in the
tail. All reported error comparisons use integer or rational arithmetic.
Floating-point quantiles only propose critical counts; they do not certify them.

\section{Calibration, timing and failure cases}
If recorded means differ from a common-model mean by at most a known absolute
amount $\delta$, use null boundaries $\epsilon_0-\delta$ and
$\beta_0+\delta$ in the two count tails. For a predicted alternative with the
same uncertainty, its adverse bounds become $p_N+\delta$ and
$p_k-\delta$. Calibration failure probability, if any, is added separately.
The supplied routine implements this conservative extension, although the
numerical design \eqref{eq:hier-errors} uses $\delta=0$.

No hidden-state preparation purity or rate bound is required. What must be
stable is the preparation distribution, repeated channel, and readout across
settings. A one-state device with a readout chosen according to $k$ can fake
any sequence $p_k$; a time-dependent or history-dependent classical process
can also invalidate the factorization. Controller state must either be reset
or included in the count. Arbitrary independent jitter at each total wait time
need not preserve the additive grid when hidden rates are unbounded; a common
random preparation latency, or independent identical increments composing each
step, does preserve the repeated-channel model. General timing uncertainty
requires a separate calibrated record-probability bound.

For context, an $N$-state deterministic discrete cycle with detector values
$\sin^2(\pi j/N)$ reproduces the ideal unitary sequence exactly. Thus the
certificate demonstrably excludes small classical models while allowing a
larger classical explanation. It is neither a proof of nonclassicality against
all models nor a new universal physical law.
